\section{\bi{Introduction}{引言}}\label{sec:introduction}
\bi{This manuscript is a candidate proof for independent review. Its mathematical conclusions and the narrower scope of selected machine checks are distinguished in Appendix~\ref{app:reproducibility}.}
{本文是供独立复核的候选证明。数学结论与选定机器检查的较窄范围在附录~\ref{app:reproducibility}中分别说明。}
\bi{The Laughlin wavefunctions provide explicit correlated states at fractional quantum Hall fillings, and contact pseudopotentials make them exact zero-energy states~\cite{Laughlin1983,Haldane1983}. A spectral-gap estimate asks a different question: can every orthogonal state be kept a positive energy above the entire zero-mode space as the particle number grows? On a torus this question also involves magnetic boundary conditions and a degenerate center-of-mass ground space. The familiar theta-function construction describes these zero modes~\cite{HaldaneRezayi1985Periodic,GreiterSchnellsThomale2016}, but their existence alone supplies no uniform lower bound on the positive spectrum.}
{Laughlin波函数在分数量子Hall填充处给出明确的关联态，而接触赝势使它们成为精确零能态~\cite{Laughlin1983,Haldane1983}。谱隙估计回答的是另一个问题：当粒子数增长时，能否让每个正交态与整个零模空间保持正的能量间隔？在环面上，这一问题还涉及磁边界条件与简并的质心基态空间。熟知的theta函数构造描述了这些零模~\cite{HaldaneRezayi1985Periodic,GreiterSchnellsThomale2016}，但零模的存在本身并不给出正谱的一致下界。}

\bi{Finite Fock-space certificates provide a useful way to separate local algebra from geometry. The fermionic certificate of OpenAI~\cite{OpenAI2026Fock} yields a uniform inequality for the full spherical interaction and its untruncated planar limit. Later work studies spherical stability~\cite{OpenAI2026Stability}, and Gao gives improved sphere-and-plane bounds through a semidefinite-program search~\cite{Gao2026ImprovedGap}. We deliberately retain the earlier seven-row certificate rather than claim its constant is the strongest available. For bosons we use the explicit 138-row certificate in the candidate proof prepared and curated by Xin Shen~\cite{Shen2026BosonicCandidate}. Only its finite planar coefficient inequalities are inputs here; neither its proposed global theorem nor its partial formalization is assumed. Those finite inequalities, and the fermionic ones, admit direct independent rational verification.}
{有限Fock空间证书把局部代数与几何分离。OpenAI的费米证书~\cite{OpenAI2026Fock}给出完整球面相互作用及其未截断平面极限的一致不等式。后续工作研究球面稳定性~\cite{OpenAI2026Stability}，Gao则通过半正定规划搜索提高球面与平面下界~\cite{Gao2026ImprovedGap}。本文有意保留较早的七行证书，不声称其常数是现有最强结果。玻色情形使用由Xin Shen整理的候选证明~\cite{Shen2026BosonicCandidate}中的明确138行证书。这里只把其有限平面系数不等式作为输入，不假设其候选整体定理或部分形式化已完成。这些有限不等式和费米不等式均可直接以独立有理算术验证。}

\bi{The geometric distinction is essential. The OpenAI source concerns spherical fermionic $V_1$ and its planar limit; Shen's candidate concerns spherical bosonic $V_0$, with a proposed sufficient flux threshold $Q\ge3088$~\cite{OpenAI2026Fock,Shen2026BosonicCandidate}. Our proposed torus thresholds concern the actual shortest lattice cycle and use $M=3N$ or $M=2N$ at Laughlin filling, in place of the spherical shifted relations. Gao's spherical lower bound $1/5$ is an energy bound at fermionic filling $1/3$, and supplies no torus gap theorem~\cite{Gao2026ImprovedGap}.}
{几何差别是必要的。OpenAI来源研究球面费米$V_1$及其平面极限；Shen的候选证明研究球面玻色$V_0$，提出充分磁通阈值$Q\ge3088$~\cite{OpenAI2026Fock,Shen2026BosonicCandidate}。本文拟议的环面阈值针对实际最短格矢，Laughlin填充关系为$M=3N$或$M=2N$，不采用球面的带移位关系。Gao的球面下界$1/5$是费米填充$1/3$处的能量下界，不提供环面谱隙定理~\cite{Gao2026ImprovedGap}。}

\bi{The torus transfer is not a replacement of spherical $SU(2)$ averaging by a formal magnetic-translation average. A finite torus pair interaction is generally not a projector, raw periodized orbitals are not orthogonal, and a finite local approximation does not control all center-of-mass modes. A three-body single-contraction identity must respect its magnetic phases and the order of the COM trace. The bosonic problem adds unbounded occupation numbers: the bounded-annihilator argument available for fermions cannot be reused. These are the specific analytic steps addressed in this manuscript.}
{环面迁移不是把球面$SU(2)$平均形式上换成磁平移平均。有限环面两体相互作用通常不是投影子，原始周期化轨道并不正交，有限局部近似也不能控制全部质心模式。三体单收缩恒等式必须正确保留磁相位与COM偏迹的次序。玻色问题还具有无界占据数，因此不能复用费米情形的有界湮灭算符论证。本文处理的正是这些具体解析步骤。}

\bi{Our common construction periodizes the planar Gaussian jets without renormalizing them. Parameter Landau-level integration then preserves relative matrix units, while a relative-jet identity controls the Gram operator on the entire output theta space. Positive finite certificates survive by congruence even when their Gram matrices have kernels. The resulting full-Fock estimates are assembled by two distinct error budgets: fermionic errors are paid for by the pair energy, whereas bosonic coherent-state comparisons pay for them jointly with the pair energy and a positive four-body term. Theorems~\ref{thm:boson} and~\ref{thm:fermion} state the resulting bounds.}
{共同构造是在不重新归一化的前提下周期化平面高斯jets。参数Landau能级积分保持相对matrix units，而相对jet恒等式控制整个输出theta空间上的Gram算符。有限正证书通过congruence保留，即使其Gram矩阵有核也如此。最终全Fock估计采用两套不同误差预算：费米误差由pair能量支付，玻色则通过相干态比较，由pair能量与正四体项共同支付。所得下界由定理~\ref{thm:boson}与~\ref{thm:fermion}陈述。}

\bi{The geometric condition is growth of the actual shortest noncontractible cycle at fixed magnetic length. Existing rigorous gap analyses for truncated thin-torus models~\cite{NachtergaeleWarzelYoung2021,WarzelYoung2022ThinTorus,WarzelYoung2023BulkEdge} concern different Hamiltonians or limiting regimes and are not used as gap inputs for the full interactions here. Pairing and coherent-state descriptions of the torus~\cite{OrtizNussinovDukelskySeidel2013,Fremling2013Coherent} inform the setting, but the all-COM frame estimate and the relative error bounds are proved below. All essential kernel, convergence, and occupation arguments are included in the manuscript. The finite exact datasets and executable checks accompany it separately from the finite-size numerical diagnostics.}
{几何条件是在固定磁长下实际最短非收缩周期增长。已有关于截断薄环面模型的严格谱隙分析~\cite{NachtergaeleWarzelYoung2021,WarzelYoung2022ThinTorus,WarzelYoung2023BulkEdge}涉及不同Hamiltonian或极限范围，不作为这里完整相互作用的谱隙输入。环面的pairing与相干态描述~\cite{OrtizNussinovDukelskySeidel2013,Fremling2013Coherent}提供背景，而全COM frame估计与相对误差界在下文证明。所有关键核、收敛及占据数论证都纳入正文和证明附录。有限精确数据及可执行核查随稿提供，并与有限尺寸数值诊断分开。}
